\section{Arithmetic fibers, scalar Hecke actions, and separation}
\label{sec:arithmetic-fibers-v43}

\status{This chapter continues the completed ordinary, control,
weight-representation, and specialization chapters.  It constructs
arithmetic points of the original weight ring, proves their density,
and develops the fiber and scalar-action arguments needed to separate
the original ordinary Hecke algebra.  The original automorphic-part
hypothesis is kept explicit wherever it is used.  Source predicates
are unpacked into the mathematical statements actually needed below;
no formal theorem is accepted in place of a written argument.}

Write $\overline{\mathbb Q}_\ell$ for the original algebraic closure,
with its original $\ell$-adic norm.  It is not replaced by its
completion.  A coefficient ring $O=c.O$ is the original finite
extension of $\mathbb Z_\ell$ inside this algebraic closure.  The
specified isomorphism $\iota:\overline{\mathbb Q}_\ell\simeq\mathbb C$
is always retained as an algebraic isomorphism.  No continuity of
$\iota$ is used.  All assertions involving an $\mathrm{HSet}$ retain
its complete original $\mathrm{Ax}$ hypothesis; in particular its
odd-prime condition includes $\ell=3$.  The weight-density argument
itself works for every prime, exactly as its source statement does.

\subsection{Dense arithmetic points in the actual inertial weight algebra}

% BEGIN AFS1
\paragraph{AFS1. Integral power series are separated by infinite Cartesian sets.}
\label{claim:AFS1}
Let $O$ be a coefficient DVR and let $S$ be a finite set.  For each
$s\in S$, let $D_s$ be an infinite subset of
$\{z\in\overline{\mathbb Q}_\ell:|z|_\ell<1\}$.  If
$f\in O[[X_s:s\in S]]$ evaluates to zero at every point of
$\prod_sD_s$, then $f=0$.  Each evaluation here is an actual
$O$-algebra map to $\overline{\mathbb Q}_\ell$.
\begin{proof}
For any finite tuple of arguments choose a finite extension $E$ of
$\mathbb Q_\ell$ containing them and the fraction field of $O$.
All arguments lie in $\mathfrak m_E$.  A series with integral
coefficients converges in $E$ at this tuple: the terms of total
degree at least $r$ have norm at most $\rho^r$, where
$\rho<1$ is the largest argument norm.  The same bound justifies
addition, multiplication, and regrouping.  Thus evaluation is the
claimed algebra map, with values in the algebraic closure itself.

We prove a stronger one-variable assertion.  If
$g(T)=\sum_{k\geq0}a_kT^k\in O_E[[T]]$ is nonzero, let
$C=\max_k|a_k|>0$ and let $N$ be the first index attaining this
maximum.  The maximum exists because the nonzero norms in a DVR
form a discrete sequence.  Then $g$ has at most $N$ distinct zeros
in the open unit disk of $\overline{\mathbb Q}_\ell$.
If $\alpha$ is one such zero, work in the finite extension
$E(\alpha)$ and put
\[
 b_j=\sum_{k>j}a_k\alpha^{k-j-1},\qquad
 h(T)=\sum_{j\geq0}b_jT^j.
\]
The sums converge, $|b_j|\leq C$, and the coefficient identities
$a_0=-\alpha b_0$, $a_j=b_{j-1}-\alpha b_j$ give
$g=(T-\alpha)h$.  If $N=0$, the term $a_0$ strictly dominates
every other term at $\alpha$, which contradicts $g(\alpha)=0$.
For $N>0$ the first coefficient of $h$ of norm $C$ is
$b_{N-1}$: its first summand is $a_N$ and all subsequent summands
have smaller norm; for $j<N-1$, the finitely many earlier terms
have norm less than $C$ and the remaining tail has norm at most
$C|\alpha|^{N-j-1}<C$.  Division therefore reduces the first
maximum index by one.  Every other zero remains a zero of $h$,
since it differs from $\alpha$.  Repeating over the finite
compositum of any proposed finite list of roots proves the bound.
This argument does not assume that an infinite root set has an
accumulation point inside the open disk.

Induct on $|S|$.  Fix all arguments but one.  Evaluation in those
arguments produces a one-variable series with coefficients in the
integer ring of one finite extension.  Its infinitely many zeros
force all its coefficients to vanish by the preceding bound.
Each coefficient is the evaluation of a series in the remaining
variables.  The induction hypothesis kills all those series.
For $S=\varnothing$ the sole evaluation is the inclusion of $O$,
so the assertion also holds.
\end{proof}
% END AFS1

% BEGIN AFS2
\paragraph{AFS2. Infinite disjoint torsion boxes with the original index condition.}
\label{claim:AFS2}
Let $n\geq1$, let $S$ be finite, and let
$\mathrm{idx}:S\to\{0,\ldots,n-1\}$.  Suppose for each $s$ that
$u_s\in\overline{\mathbb Q}_\ell$ has a positive power
satisfying $|u_s^{M_s}-1|<1$.  There are infinite sets $R_s$ of
roots of unity such that
\[
 |u_s\eta-1|<1\quad(\eta\in R_s),\qquad
 R_s\cap R_t=\varnothing\quad\hbox{if }\mathrm{idx}(s)\ne\mathrm{idx}(t).
\]
Consequently $f\in O[[X_s]]$ is zero if every algebra map $y$
with
\[
 1+y(X_s)=u_s\eta_s,
 \quad\eta_s\text{ of finite order},\qquad
 \eta_s\ne\eta_t\text{ whenever }\mathrm{idx}(s)\ne\mathrm{idx}(t)
\]
kills $f$.
\begin{proof}
The hypothesis implies $|u_s|=1$.  Put all the finitely many $u_s$
in one finite extension $E$.  If its residue field has $q$ elements,
$M=q-1$ is prime to $\ell$ and $|u_s^M-1|<1$ for every $s$.
Factoring $T^M-1$ in the algebraic closure gives a root $\xi_s$
with $|u_s-\xi_s|<1$: otherwise the product of all the distances
would be at least one.  Its norm is one, so
$\omega_s=\xi_s^{-1}$ satisfies
$\omega_s^M=1$ and $|u_s\omega_s-1|<1$.

For each $k\geq1$ choose a primitive $\ell^k$th root $\zeta_k$.
It satisfies $|\zeta_k-1|<1$: after reduction in the integer ring
of its finite extension, $T^{\ell^k}-1=(T-1)^{\ell^k}$, so its
residue is one.  Set
\[
 R_s=\{\omega_s\zeta_{\mathrm{idx}(s)+1+nt}:t\geq0\}.
\]
The $\ell$-primary part of the order of
$\omega_s\zeta_k$ is exactly $\ell^k$, because the order of
$\omega_s$ divides $M$ and $(M,\ell)=1$.  Equality of two such
elements forces equality of their $k$'s.  The displayed levels
are distinct for different $t$, and equality between levels for
$s,t$ forces their indices modulo $n$ to coincide.  This proves
infinitude and the required disjointness.  The norm inequality
follows from
$ab-1=a(b-1)+(a-1)$ for two units congruent to one.
The sets $D_s=\{u_s\eta-1:\eta\in R_s\}$ are infinite subsets
of the open unit disk.  Every evaluation on their Cartesian
product satisfies the prescribed conditions, so AFS1 applies.
\end{proof}
% END AFS2

% BEGIN AFS3
\paragraph{AFS3. Torsion on inertia coordinates gives the actual finite-part depth.}
\label{claim:AFS3}
Use the original local decomposition groups $G_v$, inertia groups
$I_v$, and stored Artin maps $a_v=\mathrm{artOf}(s.hw(v))$.
Choose the exact local abelian coordinates furnished by KX3,
\[
 a_v^{\rm ab}:G_v\longrightarrow
            \mathbb Z_\ell^{d_v}\times\mathbb Z_\ell,
 \qquad d_v=[L_{s.w(v)}:\mathbb Q_\ell]>0,
\]
and choose $\gamma_{v,j}\in I_v$ mapping to $(e_j,0)$.
If a continuous character
$\Phi:G_v\to\overline{\mathbb Q}_\ell^{\times}$ has finite-order
values on every $\gamma_{v,j}$, then $\Phi|_{I_v}$ has finite
image.  Some integer $e$ has the following property for all
$e'\geq e$: if $t$ is in the original Weil group and
$a_v(t)\in U_c(L_{s.w(v)},e')$, then $\Phi(t)=1$.
\begin{proof}
KX3 identifies the image of inertia in the topological
abelianization of the actual $G_v$ with
\[
  O_{L_{s.w(v)}}^{\times}\simeq
       T_v\times\mathbb Z_\ell^{d_v},
\]
where $T_v$ is finite, including both its prime-to-$\ell$ and
$\ell$-primary torsion.  Under this identification the selected
map $a_v^{\rm ab}$ is projection to the free coordinates on
inertia.  A continuous character into a Hausdorff abelian group
kills the closure of the commutator subgroup.  It thus factors
through this same abelianization.  Choose a common positive
multiple $N$ of the orders of its finitely many torsion values
and of the values on the $\gamma_{v,j}$.  The character $\Phi^N$
is trivial on $T_v$ and on all the coordinate generators.
Their integer span is dense in $\mathbb Z_\ell^{d_v}$, so
continuity makes $\Phi^N$ trivial on all inertia.  Its image
therefore lies in the finite group of $N$th roots of unity.
Its kernel is open, since this finite subset of a Hausdorff
field is discrete.

CIP4 identifies the stored Artin map with the same original
local reciprocity quotient, with
$\operatorname{ord}(a_v(t))=-\deg(t)$.  The character on inertia
therefore descends continuously to a finite character of the
local unit group.  An open subgroup of that group contains
$U_c(F_v,e)$ for some $e$, and then contains every
$U_c(F_v,e')$ with $e'\geq e$.  The valuation identity shows
that $a_v(t)$ is a unit precisely when $t$ belongs to inertia.
The descended character evaluated at $a_v(t)$ is $\Phi(t)$;
this proves the last assertion.  This uses the full local
abelianization already proved in KX3, rather than assuming
the private formal lemma about an arbitrary abstract
$\mathrm{LocAb}$ record.
\end{proof}
% END AFS3

% BEGIN AFS4
\paragraph{AFS4. The original algebraic character and finite-part identity.}
\label{claim:AFS4}
Fix all the literal data of
\texttt{TDenseU.isDomain\_and\_denseWt\_of\_isTheta}, including
the supplied $\mathrm{LamObj}\ r$, its characters $\kappa_{v,i}$,
the character $\theta$ satisfying $\mathrm{IsTheta}$, and any
dominant weight $\lambda$.  For every $v,i$ there is a
continuous character $\delta_{v,i}:G_v\to
\overline{\mathbb Q}_\ell^{\times}$ such that, for inertia
elements $t$ of the original Weil group,
\[
 \delta_{v,i}(t)=
  \operatorname{algPow}(s.w(v),\operatorname{lowW}(\lambda,i),a_v(t)).
\]
Choose any full witnessing character $\psi_{v,i}$ extending
$\kappa_{v,i}$.  For each original algebra point
$x:\Lambda\to\overline{\mathbb Q}_\ell$ put
\[
 \Phi_{x,v,i}=(x\circ\psi_{v,i})\,
                     \varepsilon_\ell^{\,i}\delta_{v,i}.
\]
These are continuous, and if $a_v(t)=y$ is a local unit then
\begin{equation}\label{eq:afs-finpart}
 \operatorname{finPart}(\theta,x,\lambda,v,i,y)
                       =\Phi_{x,v,i}(t)^{-1}.
\end{equation}
\begin{proof}
The finite set of continuous embeddings
$j:F_v\to\overline{\mathbb Q}_\ell$ defines on local units
the continuous character
\[
 \alpha(y)=\prod_jj(y)^{
          \operatorname{lowW}(\lambda,i)(j|_L)}.
\]
All values lie in the unit group of one finite extension:
each embedding is an isometric $\mathbb Q_\ell$-embedding
(WRP1), and finitely many such finite image fields have finite
compositum.  Extend $\alpha$ to $F_v^{\times}$ by making it
one on a chosen uniformizer.  It extends continuously to
$\widehat{F_v^{\times}}=\widehat{\mathbb Z}\times O_{F_v}^{\times}$.
Composing with the actual reciprocity identification of KX3
and CIP4 gives $\delta_{v,i}$.  Only its restriction to inertia
is required, so its unramified choice changes none of the
following identities.

ASP2 gives, for the given $x$, a number $0<\eta<1$ with
$|x(\mathfrak m_\Lambda^r)|\leq\eta^r$.  Thus composition
with $x$ carries any adically continuous character to a
continuous $\ell$-adic character: choose $r$ with
$\eta^r<\epsilon$ and use equality modulo
$\mathfrak m_\Lambda^r$ on a sufficiently small neighborhood.
The cyclotomic character and $\delta_{v,i}$ are continuous,
which proves the asserted continuity of $\Phi$.

The literal $\mathrm{IsTheta}$ identity, evaluated at $x$, is
\[
 x(\theta(e_{v,i}(y)))\,x(\kappa_{v,i}(t))\,
                     \varepsilon_\ell(t)^i=1.
\]
Here $t$ is in inertia because its Artin image is a unit,
and $\psi_{v,i}(t)=\kappa_{v,i}(t)$.  Divide this identity by
the original $\operatorname{algPow}$ factor.  The defining
formula for $\delta$ proves \eqref{eq:afs-finpart}, with the
inverse and the exponent $i$ exactly as in the source.
\end{proof}
% END AFS4

% BEGIN AFS5
\paragraph{AFS5. Every supplied original weight ring is a domain and has dense good points.}
\label{claim:AFS5}
Under precisely the quantified hypotheses of
\texttt{TDenseU.isDomain\_and\_denseWt\_of\_isTheta}, the given
$\Lambda.\Lambda$ is a domain, and for every $\lambda$ it
satisfies the original $H.\mathrm{DenseWt}(\theta,\lambda)$,
where $H=\mathrm{hset}(Q,\tau_0)$.  In expanded form,
\[
 \bigl[\forall x:\Lambda\to_O\overline{\mathbb Q}_\ell,
       \ H.\mathrm{Good}(\theta,x,\lambda)\Rightarrow x(f)=0\bigr]
                    \quad\Longrightarrow\quad f=0.
\]
\begin{proof}
KX4--KX5 give, for this supplied universal object, the unique
isomorphism
\[
 e:O[[X_{v,i,j}:v\mid\ell,\ 0\leq i<n,\ 1\leq j\leq d_v]]
                         \overset\sim\longrightarrow\Lambda
\]
and the actual finite coefficient characters $\psi^0_{v,i}$.
At the selected inertia generators its formula is
\begin{equation}\label{eq:afs-coordinate}
 \kappa_{v,i}(\gamma_{v,j})
       =\psi^0_{v,i}(\gamma_{v,j})\bigl(1+e(X_{v,i,j})\bigr).
\end{equation}
The map $e$ respects every $O$-algebra map in the original
universal property; it does not replace the prescribed $\Lambda$
by an unrelated model.  A power-series ring over a domain is a
domain: choose the least nonzero monomial in the well-ordered
graded lexicographic order; the least monomial in a product has
the product of those coefficients.  Hence $\Lambda$ is a domain.

For the fixed $\lambda$ take AFS4's $\delta$ and set
\[
 u_{v,i,j}=
 \bigl(\psi^0_{v,i}(\gamma_{v,j})
       \varepsilon_\ell(\gamma_{v,j})^i
       \delta_{v,i}(\gamma_{v,j})\bigr)^{-1}.
\]
Every factor is a unit in the integer ring of a finite extension,
so each $u_{v,i,j}$ has a positive power congruent to one.
AFS2 supplies a Cartesian family of evaluations $y$ such that
\[
 1+y(X_{v,i,j})=u_{v,i,j}\eta_{v,i,j},
\]
where all $\eta$ have finite order and their values are
different whenever their $i$ indices differ.  Put
$x=y\circ e^{-1}$.  Applying $x$ to
\eqref{eq:afs-coordinate} and canceling its unit factors gives
\begin{equation}\label{eq:afs-torsionvalues}
             \Phi_{x,v,i}(\gamma_{v,j})=\eta_{v,i,j}.
\end{equation}
AFS3 gives, for each $v,i$, a depth on which this character
is trivial on Artin preimages.  There are finitely many
$v,i$, so take one common upper bound.  Equation
\eqref{eq:afs-finpart} proves the exact $H.\mathrm{Arith}$
condition at this depth.

For $i\ne i'$ at a given $v$, choose $j=1$, which exists
because $d_v>0$ in KX3.  The original Artin image
$y=a_v(\gamma_{v,1})$ is a local unit.  By
\eqref{eq:afs-finpart} and \eqref{eq:afs-torsionvalues}, the
two finite parts at $y$ are the inverses of the distinct
$\eta_{v,i,1}$ and $\eta_{v,i',1}$.  They are distinct.
For the original $\mathrm{hset}(Q,\tau_0)$ every dominant
$\lambda$ is admitted.  This verifies all three clauses of
$H.\mathrm{Good}$, without changing the original characters.

If $f$ vanishes at every such good point, then $e^{-1}(f)$
vanishes on the entire torsion box just constructed.  AFS2
forces it to be zero, whence $f=0$.  All choices were made
after fixing $\lambda$ and before choosing the tested element
$f$; no point is selected by a coefficient of $f$.
\end{proof}
% END AFS5

% BEGIN AFS6
\paragraph{AFS6. The original increasing weight has a consistent dominant realization.}
\label{claim:AFS6}
For the original Hida setting $s$, every strictly increasing
weight $h:s.\mathrm{Wt}(n)$ has a dominant weight $\lambda$
with the literal relation $\mathrm{WtOf}(h,\lambda)$.
\begin{proof}
At a continuous embedding $\tau:F_v\to\overline{\mathbb Q}_\ell$
define
\[
 \lambda(\tau|_L)_j=h(v,n-1-j,\tau)-(n-1-j),
                      \qquad 0\leq j<n.
\]
Strict increase of the integer-valued $h$ implies
$h(v,b,\tau)-h(v,a,\tau)\geq b-a$ whenever $a\leq b$.
The displayed sequence is therefore dominant in the source's
order.  Its reversed coordinate gives
\[
 \operatorname{lowW}(\lambda,i)(\tau|_L)+i=h(v,i,\tau),
\]
which is exactly $\mathrm{WtOf}$.

This definition is consistent.  A continuous field embedding
of a local completion preserves its valuation, so its
restriction to $L$ determines the associated finite prime.
More explicitly, if distinct primes gave the same restriction,
choose integral $a,b$ with $a$ in the first prime, $b$ in the
second and $a+b=1$.  The common embedding would give both
$|a|<1$ and $|b|<1$, contradicting the ultrametric inequality
at $1$.  At the same prime, density of $L$ in its completion
and continuity show that the two embeddings coincide.  The
chosen primes $s.w(v)$ are themselves distinct because their
restrictions to the maximal real subfield are the original
$v$.  Thus each prescribed global embedding receives only
one value.  Set $\lambda=0$ at any remaining global embeddings.
The resulting dominant weight has all the required identities.
\end{proof}
% END AFS6

% BEGIN AFS7
\paragraph{AFS7. Local algebraicity and finite parts agree on the stored Artin map.}
\label{claim:AFS7}
Let $\kappa:G_v\to\overline{\mathbb Q}_\ell^{\times}$ be a
continuous character, let $w$ be an integer weight on the
continuous embeddings of $F_v$, and suppose a function $f$
on local units satisfies, for every inertia element $\sigma$,
\begin{equation}\label{eq:afs-isfin}
 \kappa(\sigma)f(a_v(\sigma))
                =\prod_\tau\tau(a_v(\sigma))^{-w(\tau)}.
\end{equation}
Then the original $\mathrm{IsLocAlg}(s.hw(v),\kappa,w)$ is
equivalent to $f$ being one on some $U_c(F_v,e)$.
For every witnessing finite-part character $\varphi$ in
$\mathrm{IsWtPart}$ one has
$\varphi(\sigma)=f(a_v(\sigma))$.
\begin{proof}
The last identity follows by cancellation of the nonzero
$\kappa(\sigma)$ from the defining $\mathrm{IsWtPart}$
identity and \eqref{eq:afs-isfin}.  The stored Artin map is
surjective from inertia to local units: its full Weil
surjectivity and $\operatorname{ord}\circ a_v=-\deg$ give
this immediately.  It is a quotient map on inertia, because
it is a continuous surjection from a compact group to the
Hausdorff compact local unit group.

If $\varphi$ has open kernel, the identity just proved shows
that it is constant on Artin fibers and descends to $f$.
Its kernel consequently gives an open subgroup of the local
units, which contains some $U_c(F_v,e)$.  Conversely, define
on inertia
\[
 \varphi(\sigma)=\kappa(\sigma)^{-1}
                       \prod_\tau\tau(a_v(\sigma))^{-w(\tau)}.
\]
It is a character, equals $f\circ a_v$, and its kernel
contains the open inverse image of $U_c(F_v,e)$.  It has
the required open kernel and the required product identity.
The original definition quantifies over every map satisfying
$\mathrm{IsLocalArtin}$.  CIP4's uniqueness identifies each
of those with the stored $a_v$, so the displayed identity
satisfies that full quantifier as well.
\end{proof}
% END AFS7

% BEGIN AFS8
\paragraph{AFS8. Exact arithmetic and good-point bridge for the original determinant ring.}
\label{claim:AFS8}
Retain all data and hypotheses of
\texttt{TWtBridge.arith\_good\_iff\_hset\_of\_theta\_incr\_weight}.
For every increasing $h$ choose $\lambda$ as in AFS6.  For
every $\lambda$ satisfying $\mathrm{WtOf}(h,\lambda)$ and
every original point $x:B.\mathrm{Pt}$, put
$x_\Lambda=x\circ\Lambda.\mathrm{toPair}(B.u)$.  Then
\begin{align*}
 x.\mathrm{Arith}(h)&\ \Longleftrightarrow\
       H.\mathrm{Arith}(\theta,x_\Lambda,\lambda),\\
 x.\mathrm{Good}(h)&\ \Longleftrightarrow\
       H.\mathrm{Good}(\theta,x_\Lambda,\lambda).
\end{align*}
The maps, inertia characters, and finite-part distinctions are
the original ones.
\begin{proof}
The universal property of $\Lambda.\mathrm{toPair}$ identifies
the image of $\kappa_{v,i}$ with the restriction of the
original $x.\psi_{v,i}$ to inertia.  The literal
$\mathrm{IsTheta}$ identity therefore gives
\[
 x.\psi_{v,i}(\sigma)\,
 x_\Lambda(\theta(e_{v,i}(a_v(\sigma))))\,
                         \varepsilon_\ell(\sigma)^i=1.
\]
Under the same reciprocity normalization, the local cyclotomic
identity is
$\varepsilon_\ell(\sigma)=\prod_\tau\tau(a_v(\sigma))$.
This is the local cyclotomic reciprocity identity on the
stored Artin map proved in CIP5.  Dividing by the
algebraic factor and using
$h(v,i,\tau)=i+\operatorname{lowW}(\lambda,i)(\tau|_L)$
gives \eqref{eq:afs-isfin} for
$\kappa=x.\psi_{v,i}$ and
$f=H.\operatorname{finPart}(\theta,x_\Lambda,\lambda,v,i)$.
AFS7 identifies local algebraicity with triviality at some
local unit depth.  A common upper bound for the finitely
many depths is exactly the single depth in $H.\mathrm{Arith}$.
The increasing-weight requirement on the determinant side
is the assumed $h.\mathrm{Incr}$.

Every finite-part inertia character is, by AFS7, exactly
$f\circ a_v$.  Since $a_v:I_v\to O_{F_v}^{\times}$ is
surjective, two such characters are equal if and only if
their functions $f$ are equal.  Thus their injective indexing
by $i$ is equivalent to the source condition
$\forall i\ne j\ \exists y,\ f_i(y)\ne f_j(y)$.
All dominant weights are admitted for the literal $H$ here.
These observations prove both equivalences, in both
directions, with the original universal-point map.
\end{proof}
% END AFS8

\subsection{Finite fibers and local scalar actions}

% BEGIN AFS9
\paragraph{AFS9. Every point of a faithful finite algebra occurs in its module fiber.}
\label{claim:AFS9}
Let $A$ be a commutative ring, let $N$ be a finite $A$-module,
and let $T\subset\operatorname{End}_A(N)$ be a commutative
$A$-subalgebra.  For a field $K$, an algebra map $x:A\to K$,
and a ring map $y:T\to K$ extending $x$, the finite-dimensional
space $V=K\otimes_A N$ contains a nonzero vector $v$ with
\[
             t v=y(t)v\qquad(t\in T).
\]
More generally, if $V=\sum_jV_j$ with every $V_j$ stable
under $T$, such a vector can be chosen in one of the $V_j$.
No flatness of $A\to K$ is assumed.
\begin{proof}
Put $\mathfrak p=\ker y$.  The module $N$ is finite over $T$
because an $A$-generating set also generates it over $T$.
It is faithful over $T$ by the specified inclusion into its
endomorphisms.  Thus $N_{\mathfrak p}\ne0$: if the localization
vanished, a product of finitely many elements outside
$\mathfrak p$ would annihilate all the chosen generators,
contradicting faithfulness.  Nakayama's lemma over the local
ring $T_{\mathfrak p}$ gives
$N_{\mathfrak p}/\mathfrak pN_{\mathfrak p}\ne0$.
The map $T/\mathfrak p\to K$ extends to its fraction field,
so extension of this nonzero vector space gives
\[
       K\otimes_{T,y}N\ne0.
\]
This is a quotient of $V$ on which $T$ acts by $y$.

Here is the finite algebra argument that turns this quotient
into an eigenvector, also keeping a stable summand.  Let
$D\subset\operatorname{End}_K(V)$ be the $K$-algebra generated
by the $T$-action.  It is commutative and finite-dimensional.
The nonzero quotient just constructed makes its scalar rule
a well-defined character $\chi:D\to K$.  A finite-dimensional
commutative algebra is a product of Artinian local algebras.
For clarity, the descending powers of its radical stabilize;
Nakayama after localization at each of its finitely many
maximal ideals makes the stable power zero.  The Chinese
remainder theorem for sufficiently large powers of the
maximal ideals then gives this product and its idempotents.
Let $e_\chi$ be the idempotent of the factor with character
$\chi$.  The space $e_\chi V$ is nonzero because it maps
onto the nonzero quotient.  Its local radical $\mathfrak r$
is nilpotent.  In a nonzero last power
$\mathfrak r^q(e_\chi V)$, every nonzero vector is killed by
$\mathfrak r$ and therefore has the required eigencharacter.

If $V=\sum_jV_j$, the scalar and $T$ stability of each $V_j$
implies $D$ stability.  Hence
$e_\chi V=\sum_je_\chi V_j$, so one $e_\chi V_j$ is nonzero.
Apply the same last-power argument within that space.
This proves both assertions, including an infinite indexing
set $j$, since membership in a sum uses finitely many terms.
The argument also shows the converse elementary fact used
below: an eigencharacter in a finite-dimensional commuting
representation occurs in a nonzero one-dimensional quotient
of its corresponding local block.
\end{proof}
% END AFS9

% BEGIN AFS10
\paragraph{AFS10. Nonzero ordinary blocks have nonzero fibers at every given point.}
\label{claim:AFS10}
For the original $H,A,\theta,a,K$ satisfying $H.\mathrm{Ax}$,
$\mathrm{IsCO}$ and $H.\mathrm{IsWt}$, set
\[
 N_0=H.X.\mathrm{Ord}(H.\tau_A,\theta,H.\mathrm{lev}(K)),
 \qquad N=H.MK(\theta,a,K),\qquad T=H.TK(\theta,a,K).
\]
If $N\ne0$, every $O$-algebra point $x:A\to
\overline{\mathbb Q}_\ell$ has a nonzero fiber
$\overline{\mathbb Q}_\ell\otimes_{A,x}N$ and extends to at
least one $O$-algebra point of $T$.  Every specified such
extension occurs as a nonzero simultaneous eigenvector in
that fiber.  At an arithmetic point, ASP's same specialization
map identifies its complex scalar extension injectively with
$\operatorname{span}_{\mathbb C}\operatorname{sp}(N)$.
\begin{proof}
OHA6--OHA7 identify $N$ with the image of an idempotent
$e\in\operatorname{End}_A(N_0)$ belonging to the original
Hecke algebra.  Both $N_0$ and $N$ are finite free over the
given local ring $A$.  Nonzero $N$ has a positive rank, so
every field fiber of it is nonzero, even for a nonflat point
map.  The image of $T$ in the endomorphisms of this fiber
generates a finite-dimensional commutative algebra over the
algebraically closed field $\overline{\mathbb Q}_\ell$.
A maximal ideal has residue field equal to that field;
the corresponding local block has a nonzero eigenvector
by the last-power argument in AFS9.  Its character restricts
to a point of $T$ extending $x$.  For any specified extension,
AFS9 applies to the actual faithful $T$-module $N$.

At an arithmetic point ASP11 says that the images under
$\operatorname{sp}$ of every $A$-basis of $N_0$ are
$\mathbb C$-linearly independent.  Therefore its tensor map
\[
 \mathbb C\otimes_{A,\iota x}N_0
     \longrightarrow\operatorname{span}_{\mathbb C}
                                      \operatorname{sp}(N_0)
\]
is an isomorphism.  The inclusion $N\hookrightarrow N_0$
is split by $e$, so it remains injective after arbitrary
scalar extension.  Restricting the preceding tensor map
gives precisely the asserted identification on $N$.
The assertion is about a nonzero specified block; it does
not assert that an arbitrary residual label produces a
nonzero block.
\end{proof}
% END AFS10

% BEGIN AFS11
\paragraph{AFS11. Stable positive-operator images inject into the original Jacquet module.}
\label{claim:AFS11}
Let $F$ be any nonarchimedean local field, let
$G=\mathrm{GL}_n(F)$ with its original $\mathrm{glIwa}$ data,
and let $\pi$ be a smooth complex representation.  Write
$N$ for the compact upper-unipotent subgroup, $N_F$ for the
full upper-unipotent group, $z=\operatorname{diag}
(\varpi^{n-1},\ldots,1)$, and $p:\pi\to J_N(\pi)$ for the
original unnormalized coinvariant map.  For $N$-fixed $v$,
\begin{equation}\label{eq:afs-jac-kernel}
 p(v)=0\quad\Longleftrightarrow\quad
                      U_z^rv=0\text{ for some }r\geq0.
\end{equation}
If $\pi$ is irreducible and the original
$\mathrm{LiteratureRankNHypothesis}$ is supplied, then for
every character $\varphi$ trivial on $T_c(c)$ the map $p$
is injective on
\[
 \operatorname{invPart}(\pi,c,\varphi)
       =\bigcap_{r\geq0}U_z^r\operatorname{eig}(\pi,c,\varphi).
\]
\begin{proof}
The integral double-coset formulas of OST5 give
$U_z^r=U_{z^r}$ on $N$-fixed vectors.  Since $p$ kills
$\pi(u)w-w$ for $u\in N_F$, the original sum over
$N/(aNa^{-1})$ gives
\begin{equation}\label{eq:afs-jac-sum}
 p(U_av)=[N:aNa^{-1}]\,p(\pi(a)v)\qquad(a\in P).
\end{equation}
These are the actual cosets and the integer index; no
normalizing measure is inserted.

If $p(v)=0$, write $v$ as a finite sum of relations
$\pi(u)w-w$ with $u\in N_F$.  Conjugation by $z^r$ carries
all these finitely many $u$ into $N$ for sufficiently large
$r$: their strictly upper entries acquire powers
$\varpi^{r(j-i)}$.  Thus $\pi(z^r)v$ belongs to the span of
$N$-relations.  Average over $N$ on the finite-dimensional
span needed for these relations.  This is a genuine finite
average: smoothness gives an open subgroup fixing the
finitely many vectors, and its normal core has finite index
in compact $N$.  The average kills every $N$-relation by
permutation of cosets.  On $\pi(z^r)v$ it equals the
nonzero scalar $[N:z^rNz^{-r}]^{-1}$ times $U_{z^r}v$,
since this vector is fixed by $z^rNz^{-r}$.  Hence
$U_z^rv=0$.  Conversely, if $U_z^rv=0$, equation
\eqref{eq:afs-jac-sum}, the invertibility of $\pi(z^r)$ on
coinvariants, and the nonzero integer index imply $p(v)=0$.

The eigenspace in the statement is fixed by $I_1(c)$:
factor an element as upper-unipotent, $T_c(c)$, and lower
depth-$c$ factors, and use its three defining eigenlaws.
The admissibility clause of the retained
$h_{\rm Lit}.\mathrm{loc.bz}$ therefore makes it finite-dimensional.
OST5--OST6 show that every $U_a$, $a\in P$, preserves it
and commutes there with $U_z$.  In a finite-dimensional
space the decreasing images of $U_z^r$ stabilize, and $U_z$
is an automorphism on their intersection.  An element of
that intersection killed by a power of $U_z$ is zero.
Equation \eqref{eq:afs-jac-kernel} proves the claimed
injectivity.  The same commutation proves that all $U_a$
preserve this intersection.
\end{proof}
% END AFS11

% BEGIN AFS12
\paragraph{AFS12. Distinct unit characters force the exact principal series and its scalars.}
\label{claim:AFS12}
Retain exactly the hypotheses of
\texttt{ExistsForm.exists\_principal\_series\_of\_\allowbreak
distinct\_inv\_ne\_bot}:
$F$ is any nonarchimedean local field, $n\geq0$, $\pi$ is a
smooth irreducible complex representation of
$\mathrm{GL}_n(F)$, $\varphi$ is a character of the original
compact diagonal torus, $\varphi(e_i(y))=1$ for
$y\in U_c(F,c+1)$, its coordinate restrictions are pairwise
distinct, and $\operatorname{invPart}(\pi,c,\varphi)\ne0$.
Assume the original $h_{\rm Lit}$.
There are smooth characters $\omega_i:F^\times\to\mathbb C^\times$
with
\[
 \varphi(e_i(y))=\omega_i(y)^{-1},\qquad
 \pi\simeq\mathrm{PS}(\omega),\qquad
 \dim_{\mathbb C}\operatorname{invPart}(\pi,c,\varphi)=1,
\]
and for every $j\leq n$ and every vector in this same
$\operatorname{invPart}$,
\begin{equation}\label{eq:afs-pseig}
 U_{a_j}v=\mathrm{psEig}(\omega,j)v,
 \qquad
 \mathrm{psEig}(\omega,j)=\prod_{0\leq i<j}
                   \omega_i(\varpi)q^{-i}.
\end{equation}
Here $q=|k_F|$, $a_j=\operatorname{diag}
(\underbrace{\varpi,\ldots,\varpi}_{j},1,\ldots,1)$,
and the source convention is
$\mathrm{PS}(\omega)=\operatorname{Ind}^{\rm norm}_B
(\omega_i|\cdot|^{(n-1)/2})_i$.
\begin{proof}
We spell out the three local literature clauses actually used.
The supplied $h_{\rm Lit}.\mathrm{loc.bz}$ says that smooth
irreducibles are admissible and that Jacquet modules of smooth
finite-length representations have finite length.  Its
$\mathrm{loc.geo}$ clause gives the normalized Jacquet
filtration of an induced representation by the actual Bruhat
cells.  Its $\mathrm{loc.disj}$ clause says that induction of
two irreducibles with totally disjoint cuspidal supports is
irreducible, where the excluded twists have exponents
$-1,0,1$.  These are explicit retained premises of this
source endpoint; the argument below does not assert their
unconditional production.

Let $J$ be the normalized Jacquet representation of $\pi$,
on the same vector space as the coinvariants in AFS11.  It
has finite length over the diagonal torus.  A smooth simple
representation of $(F^\times)^n$ is one-dimensional: an
open compact subgroup fixing one nonzero vector fixes its
entire cyclic representation, because the torus is abelian;
the resulting quotient group is finitely generated abelian.
Its complex group algebra is a finitely generated
commutative algebra, and a simple quotient is $\mathbb C$
by the weak Nullstellensatz.  Thus a finite composition
series makes $J$ finite-dimensional.  This elementary torus
argument applies also in positive residue characteristic.

The nonzero image of $\operatorname{invPart}$ under AFS11
lies in the subspace $E\subset J$ on which the compact
diagonal torus acts by $\varphi^{-1}$; the modulus character
is one on this torus.  This subspace is stable under the
full diagonal torus.  Its commuting endomorphisms have a
common eigencharacter $\mu$, and the corresponding local
block of $J$ has a nonzero quotient functional
$l:J\to\mathbb C$ of character $\mu$ (AFS9's last
finite-dimensional observation).  The character is smooth:
choose a vector on which $l$ is nonzero and an open
stabilizer of that vector.  Write
$\mu(t)=\prod_i\psi_i(t_i)$; then every $\psi_i$ is smooth
and
\[
                 \psi_i(y)=\varphi(e_i(y))^{-1}
                                      \quad(y\in O_F^\times).
\]

For completeness, the needed Frobenius map has the explicit
formula $v\mapsto[g\mapsto l(p(\pi(g)v))]$.
The definition of normalized Jacquet action gives its left
$B$-transformation law $\delta_B^{1/2}\mu$; unipotent factors
act trivially on $p$.  Smoothness follows from a stabilizer
of $v$, and $B\backslash G$ is compact, so this is the
original normalized induced representation.  The map is
nonzero and is equivariant for right translation.  By
irreducibility it embeds $\pi$ into
$\operatorname{Ind}^{\rm norm}_B(\psi)$.

The Jacquet functor is exact here.  To see the only
nontrivial direction, a relation killed in a larger module
uses finitely many upper-unipotent elements.  They lie in
one compact open subgroup of $N_F$; averaging over it as
in AFS11 shows that the same vector is already a sum of
relations in the submodule.  Applying exactness to the
preceding embedding puts $J$ inside the Jacquet module
of this principal series.  The geometric-lemma filtration
has one one-dimensional factor with character
$\prod_i\psi_{w(i)}(t_i)$ for each permutation $w\in S_n$.
Indeed the size matrices for singleton blocks have every
row and column sum one, hence are permutation matrices;
the corresponding zero- and one-dimensional block
inductions simply permute the $\psi_i$, with the normalized
modulus factors canceling.  This verifies the precise
specialization of the supplied geometric-lemma clause.

The functor taking a given smooth compact-torus character
space is exact: twist by its inverse and use the finite
average over a common finite quotient on any proposed
lift.  Its application to that filtration shows that the
$\varphi^{-1}$ subspace has dimension at most the number
of $w$ with
$\psi_{w(i)}|_{O_F^\times}=\psi_i|_{O_F^\times}$ for all $i$.
Distinctness of the given coordinate characters forces
$w=1$.  Therefore $\dim E\leq1$.  It is nonzero, and the
full torus acts on it by $\mu$.  AFS11 now makes
$\operatorname{invPart}$ one-dimensional as well.

Set $\omega_i=\psi_i|\cdot|^{-(n-1)/2}$.  The unit
restrictions remain distinct.  Induction in stages and
$h_{\rm Lit}.\mathrm{loc.disj}$ prove that the principal
series is irreducible: successively add one $\mathrm{GL}_1$
character; no character already present can equal the
new one times $|\cdot|^r$ for $r=-1,0,1$, since these
twists are trivial on units.  The cuspidal supports at
each step are precisely these singleton characters.
The starting $n=0$ representation is the one-dimensional
representation of the trivial group, and $n=1$ is a
character.  Thus the nonzero embedding above is an
isomorphism, proving the asserted principal-series statement.

Finally $[N:a_jNa_j^{-1}]=q^{j(n-j)}$, from the $j(n-j)$
upper entries crossing the two diagonal blocks.  In the
normalized Jacquet module, the torus action gives
\[
 p(\pi(a_j)v)=\delta_B(a_j)^{1/2}
                 \prod_{i<j}\psi_i(\varpi)\,p(v).
\]
Combining with \eqref{eq:afs-jac-sum} yields the scalar
\[
 q^{j(n-j)}\prod_{i<j}
     q^{-(n-1-2i)/2}\,
     \omega_i(\varpi)q^{-(n-1)/2}
   =\prod_{i<j}\omega_i(\varpi)q^{-i}.
\]
The exponent identity is
$j(n-j)-j(n-1)+j(j-1)/2=-j(j-1)/2$.
Both $U_{a_j}v$ and the stated scalar multiple of $v$
remain in $\operatorname{invPart}$ by AFS11, where $p$
is injective.  Their equal Jacquet images therefore give
\eqref{eq:afs-pseig}.  This keeps the original $\varpi$,
operator, normalization and indexing, including $j=0,n$.
\end{proof}
% END AFS12

% BEGIN AFS13
\paragraph{AFS13. Isotypic sums admit separating maps to the original irreducible.}
\label{claim:AFS13}
If $\mathrm{IsoSum}(\rho,Y,\sigma)$ holds in its literal
source sense, then $Y$ is stable under $\rho$, and every
nonzero $F\in Y$ is detected by a complex-linear map
$p:M\to\sigma.V$ satisfying
\[
 p(\rho(g)y)=\sigma.\rho(g)p(y)\quad(y\in Y),
                                  \qquad p(F)\ne0.
\]
\begin{proof}
The source predicate says that $Y$ is spanned by the images
of intertwiners $\sigma\to M$ whose ranges lie in $Y$.
Each nonzero such image is a simple subrepresentation
isomorphic to $\sigma$; this also proves stability of $Y$.
A sum of simple submodules is a direct sum of a suitable
subfamily.  To justify this, choose by the maximal principle
a maximal family whose sum is direct.  Any simple summand
not contained in that sum has zero intersection with it,
by simplicity, and could be added, a contradiction.
Consequently the chosen family spans $Y$.  The nonzero
vector $F$ has a nonzero component in one of its summands.
Projection to that summand followed by an isomorphism to
$\sigma$ detects $F$ and intertwines on $Y$.  Extend this
linear map from $Y$ to $M$ using a vector-space complement.
Only intertwining on $Y$ is asserted, exactly as in the
source's $\mathrm{isoSep}$ statement.
\end{proof}
% END AFS13

% BEGIN AFS14
\paragraph{AFS14. Scalar torus action and surjective positive operators on a local part.}
\label{claim:AFS14}
Let $Y$ be an isotypic sum of a smooth irreducible
$\sigma$ for $\mathrm{GL}_n(F)$ in a complex representation
$\rho$.  Let $0\ne E\subset Y$ be fixed by $N$ and
$N^-(e)$, with
\[
 \rho(e_i(y))F=f_i(y)^{-1}F\quad(F\in E),
 \qquad f_i|_{U_c(F,e+1)}=1,
\]
and suppose the $f_i$ are pairwise distinct as functions
on local units.  Suppose each $U_{a_{j+1}}$ maps $E$ onto
$E$ for $0\leq j<n$.  Under the original $h_{\rm Lit}$,
there are smooth characters $\omega_i$ such that
\[
 \sigma\simeq\mathrm{PS}(\omega),\qquad
 \omega_i(y)f_i(y)=1,
 \qquad U_{a_j}F=\mathrm{psEig}(\omega,j)F
                         \quad(F\in E,\ j\leq n).
\]
\begin{proof}
Every compact diagonal element is a product of the
single-coordinate $e_i(y)$.  Thus it acts as a scalar on
$E$.  Since $E\ne0$, that scalar is unique; it is nonzero
because the group action is invertible.  Comparing its
action on one fixed nonzero vector proves multiplicativity.
The inverse scalars define a character $\varphi$ with
$\varphi(e_i(y))=f_i(y)$.  The specified depth makes
$\varphi$ trivial on $T_c(e)$.

The original identity $z=\prod_{j=0}^{n-1}a_j$ and OST5
show that $U_z$ maps $E$ onto itself.  For any map $p$ in
AFS13, the image of $E$ satisfies the defining eigenlaws
for $\operatorname{eig}(\sigma,e,\varphi)$.  Repeatedly
taking a preimage in $E$ under $U_z$ and using naturality
of the finite sum gives
$p(E)\subset\operatorname{invPart}(\sigma,e,\varphi)$.
AFS13 supplies one $p$ with nonzero image, so AFS12 applies
and gives $\omega$ and its scalar laws on this
$\operatorname{invPart}$.

Fix $F\in E$.  The difference
$U_{a_j}F-\mathrm{psEig}(\omega,j)F$ belongs to $Y$,
since $Y$ is stable under the local group and hence under
the finite sum.  Every AFS13 map kills it, by naturality
and AFS12's scalar identity on $p(E)$.  The separation
in AFS13 forces the difference to be zero.  The assertion
uses all intertwiners only after the single $\omega$ has
been chosen from one nonzero image; their local
$\operatorname{invPart}$ is the same, so no incompatible
character choices occur.
\end{proof}
% END AFS14

% BEGIN AFS15
\paragraph{AFS15. Surjectivity forces integral eigenvalues to be units.}
\label{claim:AFS15}
Let $N$ be a finite module over a commutative ring $A$ and
let $U$ be a surjective endomorphism.  Then
\[
                  1+Uq(U)=0\quad\text{on }N
\]
for some $q\in A[T]$.  If a semilinear map $s$ carries $U$
to a complex operator $T$, its scalar ring map is
$\iota x$, every $|x(a)|\leq1$, and a nonzero eigenvector
of $T$ in $\mathbb C\operatorname{span}s(N)$ has eigenvalue
$\mu$ with $|\iota^{-1}(\mu)|\leq1$, then
$|\iota^{-1}(\mu)|=1$.
\begin{proof}
Choose generators $f_1,\ldots,f_r$ and, by surjectivity,
write $f_i=U\sum_ja_{ij}f_j$.  The adjugate identity for
$I-T(a_{ij})$ gives a polynomial annihilating these
generators when $T=U$.  Its constant coefficient is one,
so it is $1+Tq(T)$.  The identity transfers through $s$
and then to its complex span, because sums, products and
scalar multiplication commute with the semilinear map.
On the nonzero eigenvector it gives
\[
  1+\widetilde\mu\,q_x(\widetilde\mu)=0,
              \qquad\widetilde\mu=\iota^{-1}(\mu).
\]
Every coefficient of $q_x$ has norm at most one, and so
does $\widetilde\mu$.  Hence $|q_x(\widetilde\mu)|\leq1$.
Taking norms in the identity gives
$1\leq|\widetilde\mu|$, proving equality.  This determinant
argument requires neither freeness of $N$ nor flatness
of the coefficient map.
\end{proof}
% END AFS15

% BEGIN AFS16
\paragraph{AFS16. Original Hecke operators preserve every original part.}
\label{claim:AFS16}
For every original $H$, fixed $\iota,\lambda,\xi$, and
$Y:H.\mathrm{PartAltAlt}(\iota,\lambda,\xi)$, each
$H.\operatorname{opC}(\iota,\lambda,d,i)$ preserves $Y.Y$.
On vectors fixed by the product $N$, the positive operator
supported at a place $v$ is exactly the local $U_{a_j}$
at $v$, with the same finite sum and index.
\begin{proof}
AFS13 applied to each $Y.\mathrm{isoL}(v)$ shows that
$Y.Y$ is stable under every single local factor.  Any
element of the finite product is a product of these
single-factor elements, so the full product action
preserves $Y.Y$.  A positive Hecke operator is a finite
sum of these group actions, multiplied by the prescribed
$\operatorname{lowVal}^{-1}$; hence it preserves $Y.Y$.
A spherical generator acts as a scalar on all of $Y.Y$
by the literal $Y.\mathrm{sph}$ field.  This proves the
first assertion for both summands of $H.\mathrm{Op}$.

For the second, take representatives $R$ for the local
quotient $N_v/a_jN_va_j^{-1}$.  Insert each representative
in the $v$th factor and put identity in every other factor.
They represent the product quotient, because the conjugate
subgroup is unchanged at the other factors.  The original
finite sum is thus term by term the local sum on any
product-$N$-fixed vector.  This is precisely the source's
\texttt{uOp\_factor} identity; no normalization or
replacement uniformizer has been introduced.
\end{proof}
% END AFS16

% BEGIN AFS17
\paragraph{AFS17. Full original good-scalar endpoint.}
\label{claim:AFS17}
Retain the complete original hypotheses of
\texttt{goodScalB\_alt}: $h_{\rm Lit}$, $H.\mathrm{Ax}$,
$A$ of class $\mathrm{IsCO}$, $H.\mathrm{IsWt}(\theta)$,
the fixed $\iota$, a given point $x:A\to_O
\overline{\mathbb Q}_\ell$, and any supplied original
$\mathrm{IsSpec}(\theta,\iota,x,\lambda,\xi,c_0,\operatorname{sp})$.
Suppose $H.\mathrm{Good}(\theta,x,\lambda)$, fix an original
part $Y$, level $K$, and $e\geq c_0$, and assume
\[
 E=Y.Y\cap H.\mathrm{OrdAt}(\theta,\iota,x,\lambda,\xi,K,e)
                                                        \ne0.
\]
There are smooth characters $\omega_{v,i}:F_v^\times\to
\mathbb C^\times$ such that
\begin{gather*}
 Y.\sigma_l(v)\simeq\mathrm{PS}(\omega_v),\qquad
 \omega_{v,i}(y)\,
       \iota(\operatorname{finPart}_{v,i}(y))=1,\\
 \mu_{v,j}=\operatorname{lowVal}(\iota,\lambda,v,j+1)^{-1}
                         \mathrm{psEig}(\omega_v,j+1),\qquad
 |\iota^{-1}(\mu_{v,j})|_\ell=1,\\
 H.\operatorname{opC}(\iota,\lambda,d,(v,j))F
                   =\mu_{v,j}F\quad(F\in E, 0\leq j<n).
\end{gather*}
\begin{proof}
Put $S=H.\mathrm{OrdAt}(\theta,\iota,x,\lambda,\xi,K,e)$.
It is finite-dimensional because it lies in the classical
space, whose finiteness is an $\mathrm{IsSpec}$ field.
Its span field identifies $S$ with the complex span of
$\operatorname{sp}(N_0)$ for the original full ordinary
module $N_0$ of AFS10.  The exact ordinary theorem OST25
makes every positive $U_a$ an automorphism of $N_0$.
The original specialization operator identity therefore
makes every normalized local positive operator $T_{v,j}$
map $S$ onto $S$.  In particular it is invertible there.

AFS16 makes $Y.Y$ stable under $T_{v,j}$.  Thus its
restriction to $E=Y.Y\cap S$ is injective.  Since $E$ is
finite-dimensional, it is surjective onto $E$ as well.
The factor $\operatorname{lowVal}(v,j+1)$ is nonzero,
being a product of field images of a uniformizer.
AFS16's product-to-local identity now says that the
unnormalized local $U_{a_{j+1}}$ maps $E$ onto itself.

Membership in the literal classical space gives, at each
$v$, the $N_v$ and $N_v^-(e)$ fixedness and the equations
\[
 \rho_v(e_i(y))F=
             \iota(\operatorname{finPart}_{v,i}(y))^{-1}F.
\]
The supplied specialization is trivial on $U_c(F_v,c_0+1)$;
since $e\geq c_0$, it is also trivial on $U_c(F_v,e+1)$.
The good-point condition says that these $i$-indexed
finite parts are pairwise distinct.  The fixed $\iota$
preserves that distinction.  Apply AFS14 with the literal
$Y.\mathrm{isoL}(v)$ at each place.  It supplies $\omega_v$,
the principal-series isomorphism, and the scalar identity.
Multiplication by the unchanged $\operatorname{lowVal}^{-1}$
gives the displayed $\mu_{v,j}$.

Choose any nonzero $F_0\in E$.  The integrality field of
the given $\mathrm{IsSpec}$ shows
$|\iota^{-1}(\mu_{v,j})|\leq1$, because $F_0$ lies in
the original classical space at the same $K,e$.
Apply AFS15 to $N_0$, its surjective original local
operator, and the same specialization map.  ASP2 gives
$|x(a)|\leq1$ for every $a\in A$.  Thus the norm is one.
The proof has preserved every $K,e\geq c_0$, every
supplied specialization datum, and the source index
$j+1$; it does not impose a domain or flatness condition
on $A$.
\end{proof}
% END AFS17

\subsection{The original action on arithmetic fibers and separation}

% BEGIN AFS18
\paragraph{AFS18. Tensor transport and projection preserve the actual part decomposition.}
\label{claim:AFS18}
Let $N=eN_0$ be an idempotent summand of a finite free
$A$-module, where $e$ belongs to the algebra generated
by specified endomorphisms $U_i$ of $N_0$.  Let a map
$s:N_0\to W$ be semilinear for $\chi:A\to K$, with
$K$ a field, and suppose the images of one $A$-basis of
$N_0$ are $K$-linearly independent.  Suppose
$sU_i=T_is$, where $T_i$ are endomorphisms of $W$.
Set $O'=K\operatorname{span}s(N_0)$ and
$E=K\operatorname{span}s(N)$.  If
\[
 O'=\sum_j(Z_j\cap O')
\]
and every $Z_j$ is stable under all $T_i$, then
\[
 E=\sum_j(Z_j\cap E).
\]
Moreover there is one ring action
\[
 \operatorname{act}:T\longrightarrow\operatorname{End}_K(E),
 \quad T=A[U_i|_N]\subset\operatorname{End}_A(N),
\]
which carries $s(f)$ to $s(tf)$, carries each generator
to $T_i|_E$, and sends $a\in A$ to multiplication by
$\chi(a)$.  Its construction is valid for arbitrary,
possibly nonflat $\chi$.
\begin{proof}
Independence gives the tensor isomorphism
$\Psi_0:K\otimes_A N_0\simeq O'$ by
$k\otimes f\mapsto ks(f)$.  The splitting of $N\subset N_0$
survives every tensor product, so its restriction gives
$\Psi:K\otimes_A N\simeq E$.  Define
\[
 \operatorname{act}(t)=\Psi(1\otimes t)\Psi^{-1}.
\]
It is a ring action, and its stated identities follow
on $s(N)$ and hence on their span.  The same construction
on $N_0$ gives the endomorphism $e_K$ of $O'$.  It fixes
$E$, has image $E$, and is idempotent, by the original
identities for $e$ before base change.

An element of the algebra generated by the $U_i$
transports to a finite polynomial expression in the
$T_i$ and the scalars $\chi(a)$.  It preserves each
$Z_j\cap O'$.  In particular, applying $e_K$ to a finite
expression $F=\sum_jF_j$ with $F_j\in Z_j\cap O'$ and
$F\in E$ gives
$F=\sum_je_KF_j$, with $e_KF_j\in Z_j\cap E$.
This proves the asserted decomposition.  Tensor
injectivity was obtained from the split inclusion,
not inferred from a flatness assumption on $K$.
\end{proof}
% END AFS18

% BEGIN AFS19
\paragraph{AFS19. Full original part-action endpoint on each arithmetic fiber.}
\label{claim:AFS19}
Retain every hypothesis of \texttt{partAct0\_alt}, including
the original $H.\mathrm{AutomorphicHeckeSet}(\iota,P)$,
$\lambda\in H.\mathrm{admitted}$, any supplied
$\mathrm{IsSpec}$ datum at $x$, any residual label $a$,
any level $K$, and the stated commutativity on $H.MK$.
Write $N_0,N,T$ as in AFS10 and
$E=\operatorname{span}_{\mathbb C}\operatorname{sp}(N)$.
Then
\[
 E=\sum_{Y:\,P(\lambda,Y.r)}(Y.Y\cap E),
\]
which in particular proves the source's disjunction of
this equality with $\neg H.\mathrm{Good}(\theta,x,\lambda)$.
There is a single ring action
$\operatorname{act}:T\to\operatorname{End}_{\mathbb C}(E)$
with all the following properties:
\begin{enumerate}
\item For every $t\in T$ and $f\in N$,
$\operatorname{act}(t)\operatorname{sp}(f)=\operatorname{sp}(tf)$.
\item It sends every original generator to the restriction
of the same $H.\operatorname{opC}$, and sends $b\in A$
to multiplication by $\iota(x(b))$.
\item Every original $O$-algebra point $y:T\to
\overline{\mathbb Q}_\ell$ extending $x$ has a nonzero
simultaneous eigenvector $F$ in
$Y.Y\cap H.\mathrm{OrdAt}(\theta,\iota,x,\lambda,\xi,K,c_0)$
for some original part $Y$ with $P(\lambda,Y.r)$.  Its
generator eigenvalues are exactly $\iota(y(\mathrm{gen}_i))$.
\item If $x$ is good and $Y$ is any original part with
$Y.Y\cap E\ne0$, there is one $O$-algebra point
$y:T\to\overline{\mathbb Q}_\ell$ extending $x$ such
that, for every $t\in T$ and every $F\in Y.Y\cap E$,
$\operatorname{act}(t)F=\iota(y(t))F$.
\end{enumerate}
\begin{proof}
OHA6--OHA7 supply the original idempotent $e$ with
$N=eN_0$, its membership in the full ordinary Hecke
algebra, and finite freeness.  ASP's basis statement
(or the corresponding field of the supplied
$\mathrm{IsSpec}$) gives the tensor injectivity required
in AFS18.  Its operator field gives the exact
intertwining of every generator.  Its span field says
\[
 O'=\operatorname{span}_{\mathbb C}\operatorname{sp}(N_0)
         =H.\mathrm{OrdAt}(\theta,\iota,x,\lambda,\xi,K,c_0).
\]
CFD12, applied to the supplied original
$\mathrm{AutomorphicHeckeSet}$, decomposes $O'$ as the
sum of the $Y.Y\cap O'$ for precisely the parts with
$P(\lambda,Y.r)$.  AFS16 gives their generator stability.
AFS18 now gives the displayed decomposition of $E$ and
one action with the first two properties.

For a specified point $y$, AFS9 applies to the faithful
finite $T$-module $N$ and the character $\iota\circ y$.
Transport its fiber through $\Psi$ from AFS18.  The
second assertion of AFS9 uses the just-proved decomposition
to place a nonzero eigenvector in one $Y.Y\cap E$.
Since $E\subset O'$, it lies in the required ordinary
space at depth $c_0$.  This gives property 3, including
the source's existential depth by choosing exactly $c_0$.

For property 4, fix $Y$ with $Z=Y.Y\cap E\ne0$.
Spherical generators act as scalars on all of $Y.Y$.
AFS17, applied to $Y.Y\cap O'\ne0$, makes every local
positive generator scalar on $Y.Y\cap O'$ and thus on
$Z$.  Consequently every element of $T$, being a
finite algebra expression in these generators and
$A$, acts as a scalar on $Z$.  Choose once a nonzero
$F_0\in Z$.  A scalar acting on all $Z$ is unique, since
its action on $F_0$ determines it.  Denote that scalar
by $z(t)$.  The identities for $0,1,+,\cdot$ in the
ring action prove that $z:T\to\mathbb C$ is a ring
homomorphism, and the coefficient identity proves
$z|_A=\iota\circ x$.  Put $y=\iota^{-1}\circ z$.
The original $O$-algebra structures commute because
they do so on $A$.  Thus $y$ is exactly the required
$O$-algebra point, chosen before all $t$ and $F$ in
property 4.  This verifies every clause of the literal
source endpoint, including parts not in the selected
predicate family in its last clause.
\end{proof}
% END AFS19

% BEGIN AFS20
\paragraph{AFS20. Original good arithmetic points separate the original Hecke algebra.}
\label{claim:AFS20}
Under every original hypothesis of
\texttt{TCoreSep.eq\_zero\_of\_vanish\_at\_good\_points\_dense\_weight},
fix $\lambda\in H.\mathrm{admitted}$ with
$H.\mathrm{DenseWt}(\theta,\lambda)$.  For every $t\in
T=H.TK(\theta,a,K)$,
\[
 \left[
 \forall y:T\to_O\overline{\mathbb Q}_\ell,
 \ H.\mathrm{Good}(\theta,y|_A,\lambda)\Rightarrow y(t)=0
 \right]\quad\Longrightarrow\quad t=0.
\]
This retains an arbitrary original $A$ of class
$\mathrm{IsCO}$ and the supplied automorphic-part hypothesis.
It does not presume
reduced fibers or a flat point map.
\begin{proof}
Choose an $A$-basis $(b_i)$ of the original full ordinary
module $N_0$, using OHA7.  To show $t=0$ as its specified
endomorphism of $N=H.MK$, fix $f\in N$ and put $g=tf$,
viewed in $N_0$.  Fix any good point $x:A\to_O
\overline{\mathbb Q}_\ell$.  ASP13 supplies the exact
specialization datum at this $x$ and $\lambda$.  Apply
AFS19 to its space $E$.

Each nonzero piece $Y.Y\cap E$ has by AFS19(4) one
point $y$ extending $x$ on which the entire $T$-action
is scalar.  The hypothesis gives $y(t)=0$, so
$\operatorname{act}(t)$ kills that piece.  It also kills
zero pieces.  The decomposition in AFS19 therefore
shows that it kills all of $E$.  Its compatibility on
$\operatorname{sp}(f)$ gives $\operatorname{sp}(g)=0$.

Write $g=\sum_i a_i b_i$ in the chosen basis.  This is
a finite sum.  Additivity and semilinearity of the same
specialization map give
\[
             0=\sum_i\iota(x(a_i))\operatorname{sp}(b_i).
\]
Its basis-independence clause implies $x(a_i)=0$ for
every $i$.  The point $x$ was arbitrary among all good
points at the fixed $\lambda$.  DenseWt makes every
$a_i$ zero, hence $g=0$.  Since $f$ was arbitrary,
the original endomorphism $t$ is zero.

In particular this proves the exact $H.\mathrm{Out.sep}$
field without requiring the other $\mathrm{Out}$ fields.
For example it also implies that $T$ is reduced under
these hypotheses: a nilpotent element has zero image
under every field point, and the result just proved
annihilates it.  This reducedness is a consequence of
the scalar-part argument, not an input to it.
\end{proof}
% END AFS20

% BEGIN AFS21
\paragraph{AFS21. Separation descends to the very same kernel ideal.}
\label{claim:AFS21}
Let $O,A,B,T,E$ be commutative rings, with their stated
$O$-algebra structures, let $T$ also be an $A$-algebra,
and let
\[
 \lambda:A\to_O B,\qquad f:B\to_O T,
          \qquad f\lambda=\operatorname{struct}_{A,T}.
\]
Let $P$ be a predicate on ring maps $A\to E$ and $G$ a
predicate on $O$-algebra maps $B\to E$.  Assume
$P(x\lambda)\Rightarrow G(x)$ for every such $x$, and
assume the $O$-algebra points $z:T\to E$ with
$P(z|_A)$ separate $T$.  If $b\in B$ vanishes at
every point $x:B\to_O E$ with
$\ker f\subset\ker x$ and $G(x)$, then $b\in\ker f$.
\begin{proof}
For any separating point $z:T\to_O E$, let $x=z\circ f$.
It kills $\ker f$, and
$x\lambda=z\circ\operatorname{struct}_{A,T}$ has property
$P$.  Thus $G(x)$ holds and the hypothesis on $b$ gives
$z(f(b))=x(b)=0$.  Separation in $T$ yields $f(b)=0$.
This is exactly $b\in\ker f$.  Neither surjectivity of
$f$ nor flatness of any map was used.  The conclusion
keeps the literal ideal $\ker f$, without passing to
its radical, a larger support ideal, or a different
quotient.
\end{proof}
% END AFS21

% BEGIN AFS22
\paragraph{AFS22. Exact density on an original IsJ ideal, with the remaining automorphic supplier exposed.}
\label{claim:AFS22}
Fix all the original data of \texttt{TJDense}, including
$h_P:\mathrm{HypothesesH'}$, the prime $\ell$, CM field
$L$, algebraic closure, Hida setting $s$, $n\geq1$,
$V,Q,c,r,\iota$, the supplied original determinant ring
$B$, weight ring $\Lambda$, and
$R:\mathrm{UnitaryHeckeData}(Q,\iota,\Lambda)$.
Retain $V.S'=s.S$ and any original auxiliary level $K$.
The following additional installation input is kept
explicit:
\begin{equation}\label{eq:afs-aut-input}
 H.\mathrm{AutomorphicHeckeSet}
       (\iota,\mathrm{UnitaryPointProperties}(r,\iota)),
             \qquad H=\mathrm{hset}(Q,R.\tau_0).
\end{equation}
Then for every original ideal $J\subset B.B$,
\[
            R.\mathrm{IsJ}(B,K,J)\quad\Longrightarrow\quad
                                      B.\mathrm{Dense}(J).
\]
In full quantified form, for every increasing original
weight $h$ and every $b\in B.B$,
\[
 \bigl[\forall x:B.\mathrm{Pt},\ J\subset\ker x
              \ \land\ x.\mathrm{Good}(h)\Rightarrow x(b)=0\bigr]
                                      \quad\Longrightarrow\quad b\in J.
\]
\begin{proof}
The literal $R.\mathrm{IsJ}$ has two cases.  In its
zero-module case $J=\top$, so the conclusion is immediate.
In its other case it supplies the original commutativity
and coefficient structures, a surjective $O$-algebra map
\begin{gather*}
 f:B.B\longrightarrow T=H.TK(R.\theta,R.a,K),\qquad J=\ker f,\\
 f\circ\Lambda.\mathrm{toPair}(B.u)
                          =\operatorname{struct}_{\Lambda,T},
\end{gather*}
together with the stated universal determinant and
character compatibility.  These maps and this ideal
are retained exactly; the density argument needs only
the displayed compatibility from that supplied record.

Fix $h$ with $h.\mathrm{Incr}$.  AFS6--AFS8 give one
$\lambda$ with $\mathrm{WtOf}(h,\lambda)$ and the exact
equivalence between $x.\mathrm{Good}(h)$ and goodness
of its original weight restriction.  AFS5 gives
$H.\mathrm{DenseWt}(R.\theta,\lambda)$.
FNL16 supplies $H.\mathrm{Ax}$, and the supplied $R$
already has its literal $\mathrm{IsWt}$ field.
Apply AFS20 using \eqref{eq:afs-aut-input}.  It gives
separation of this exact $T$ by the good points at
$\lambda$.  Now apply AFS21 with
$A=\Lambda.\Lambda$, $\lambda=\Lambda.\mathrm{toPair}(B.u)$,
$P$ the HSet good-point predicate, and $G$ the original
$x.\mathrm{Good}(h)$.  AFS8 verifies the required
predicate implication.  The conclusion is
$b\in\ker f=J$, as required, for this same $J$ and
every $h,b$ in the source's $\mathrm{Dense}$ definition.

The source proves \eqref{eq:afs-aut-input} from the
type field of $R$ through its actual automorphic-part
and Galois-attachment constructions.  The present
chapter does not replace those producers by the
existence of a formal proof.  Thus AFS22 proves the
complete same-ideal density implication with this
exact remaining supplier, while the unconditional
installation of the bare \texttt{TJDense} endpoint
remains open until that supplier is reconstructed.
It neither asserts the other $\mathrm{Out.pt}$ fields
nor closes the complete Hida construction.
\end{proof}
% END AFS22
